% Удиртгал

\phantomsection
\addchaptertocentry{Удиртгал}
\label{Introduction}

\section*{Удиртгал}

Ясны хугарал бол дэлхий даяар хамгийн өргөн тархсан гэмтлийн нэг бөгөөд жил бүр 178 сая гаруй шинэ тохиолдол бүртгэгддэг. Монгол Улсад ч мөн 2019 оны статистикаар 156,723 осол гэмтлийн тохиолдол бүртгэгдсэн бөгөөд хөдөөд радиологичдын хомсдол нь цаг алдалгүй оношилгооны явуулахад ноцтой бэрхшээл болдог. Уламжлалт рентген зургийн оношилгоо нь хүний нүдний үзлэгт тулгуурладаг тул алдааны хувь 3--9\%-д хүрч, ялангуяа нарийн болон бүрэн бус шилжилтгүй хугарлуудыг алдах эрсдэл өндөр байдаг.

Сүүлийн жилүүдэд гүн сургалт (deep learning) нь эмнэлгийн дүрс оношилгоонд нэвтрэн орсноор энэ асуудлыг хэсэгчлэн шийдэх боломж нээгдсэн. Гэвч одоо хүртэл ихэнх судалгаанууд нь хугарал илрүүлэх, сегментлэх, хугарлын зэрэг тодорхойлох зэрэг даалгавруудыг тусад нь шийдэж байгаа бөгөөд эдгээрийг нэгдсэн дамжлагад нэгтгэж, эдгэрэлтийн хугацааг ч мөн таамаглах оролдлого хомс байна.

Энэхүү дипломын ажил нь рентген зургаас ясны хугарлыг автоматаар илрүүлж, байршлыг пиксел түвшинд тогтоон, хугарлын зэргийг ангилан, эдгэрэлтийн хугацааг таамаглах олон даалгаварт гүн сургалтын систем боловсруулахад чиглэсэн. Архитектурын суурь нь ImageNet дээр урьдчилан сургагдсан EfficientNet-B3 хуваалцсан кодлогч (shared encoder) бөгөөд дөрвөн даалгаварт модуль нэгдсэн байна. Туршилтад MURA (40,000 зураг), FracAtlas (4,083 зураг), GRAZPEDWRI-DX (20,327 зураг) зэрэг нийтэд нээлттэй өгөгдлийн сангуудыг ашигласан.

Ажил нь долоон бүлгээс бүрдэнэ. Бүлэг~\ref{Chapter1}-д судалгааны үндэслэл, зорилго, зорилт, хамрах хүрээ, таамаглалыг тодорхойлно. Бүлэг~\ref{Chapter2}-т холбогдох онолын суурь болон ижил төстэй судалгаануудын тойм өгнө. Бүлэг~\ref{Chapter3}-т судалгааны арга зүй болон санал болгох архитектурыг нарийвчлан тайлбарлана. Бүлэг~\ref{Chapter4}-т өгөгдөл цуглуулах, боловсруулах, хуваах үйл явцыг дүрслэнэ. Бүлэг~\ref{Chapter5}-д загварын хэрэгжүүлэлт болон сургалтын дамжлагыг тайлбарлана. Бүлэг~\ref{Chapter6}-д туршилтын үр дүнг танилцуулан хэлэлцэнэ. Бүлэг~\ref{Chapter7}-д эмнэлгийн хиймэл оюун ухааны ёс зүйн асуудлыг авч үзнэ.
